\chapter*{CONCLUSION}
\addcontentsline{toc}{chapter}{Conclusion}
\markboth{CONCLUSION}{}

Ce projet de six mois au sein de l'équipe Advanced Data \& Climate d'Allianz France a permis de mener à bien la première migration opérationnelle d'un datamart SAS vers une architecture Python/PySpark en contexte industriel. Conçu comme un PoC pilote du programme SAS EXIST, ce projet dépasse largement le simple exercice technique : il constitue à la fois une démonstration de faisabilité, une méthodologie formalisée et un ensemble de livrables directement réutilisables pour les migrations à venir.

Sur le plan technique, l'objectif principal (garantir la parité fonctionnelle avec le système SAS) a été atteint pour les pipelines PTF Mouvements et Capitaux. Les comparaisons réalisées sur la vision de référence montrent un alignement strict : mêmes volumes de lignes, mêmes répartitions par pôle et par produit, même structure de colonnes. La démarche itérative adoptée, fondée sur cinq phases clairement définies, une validation continue et une résolution des divergences par investigation systématique, a permis d'identifier et de corriger des différences parfois subtiles qui n'auraient pas été détectées par une simple comparaison des totaux. La migration ne s'est donc pas limitée à « traduire » du SAS vers Python : elle a permis de réarchitecturer entièrement le datamart selon les principes de l'architecture médaillon (Bronze/Silver/Gold), en externalisant la configuration et en produisant un code modulaire pensé pour durer.

Sur le plan méthodologique, ce projet a mis en évidence les conditions réellement nécessaires à une migration de ce type. L'absence de documentation du code SAS d'origine a constitué le premier obstacle : comprendre et reconstituer les règles métier à partir de macros imbriquées a représenté une part importante de l'effort global. Ce constat souligne l'importance de la documentation fonctionnelle produite dans ce PoC, qui servira de base pour les prochaines migrations. De même, la disponibilité des données de test en environnement Azure s'est révélée être un facteur critique : les protocoles de sécurité stricts ont entraîné des délais non anticipés, rappelant qu'une migration en contexte industriel implique nécessairement des contraintes opérationnelles fortes.

Pour Allianz France, les bénéfices attendus s'articulent autour de trois axes. Le premier est économique, avec la perspective de réduire progressivement les coûts liés aux licences SAS. Le deuxième est technique : le traitement distribué sur ADP offre une scalabilité bien supérieure à celle de l'environnement SAS historique. Enfin, la capitalisation documentaire produite dans le cadre de ce projet constitue un actif durable, en rendant explicite une connaissance métier qui était jusque-là intégralement noyée dans quinze mille lignes de code SAS non maintenu.

Les perspectives ouvertes par ce PoC sont immédiates. Le programme SAS EXIST prévoit la migration d'autres marchés ABR (notamment Auto et Habitation). Le framework développé (lecteur générique, infrastructure de logging, configuration externalisée, structure de projet et méthodologie de validation) est conçu pour être réutilisé sans modification majeure. La prochaine étape naturelle consiste à finaliser la validation de parité sur dix visions historiques, puis à engager la migration du deuxième datamart en capitalisant sur l'ensemble des enseignements de ce premier cycle.

Sur le plan personnel et professionnel, ce projet a constitué une opportunité précieuse de confronter les connaissances académiques aux exigences d'un environnement industriel de grande envergure. Il a favorisé une forte montée en compétences techniques, notamment sur les architectures Big Data (PySpark, Delta Lake) et les environnements Cloud (Allianz Data Platform). Par ailleurs, l'immersion au sein d'une équipe transversale a permis de mieux appréhender les enjeux stratégiques liés à la valorisation et à la gouvernance de la donnée dans le secteur assurantiel. Cette expérience formatrice confirme un intérêt marqué pour les métiers du Data Engineering et pose des bases solides pour la future carrière professionnelle.